\documentclass[12pt,a4paper]{report}

\usepackage[utf8]{inputenc}
\usepackage[T1]{fontenc}
\usepackage[spanish]{babel}
\usepackage{listings}
\usepackage{xcolor}
\usepackage{geometry}
\usepackage{hyperref}
\usepackage{parskip}
\usepackage{titlesec}
\usepackage{enumitem}
\usepackage{booktabs}
\usepackage{tcolorbox}
\usepackage{amsmath}
\usepackage{array}

\geometry{margin=2.5cm}

% ─── Colores ──────────────────────────────────────────────────────────────────
\definecolor{lisp-bg}{RGB}{248,248,252}
\definecolor{lisp-kw}{RGB}{0,0,180}
\definecolor{lisp-mac}{RGB}{0,120,0}
\definecolor{lisp-str}{RGB}{160,0,0}
\definecolor{lisp-cmt}{RGB}{110,110,110}
\definecolor{box-blue}{RGB}{240,248,255}
\definecolor{box-blue-border}{RGB}{100,149,237}
\definecolor{box-green}{RGB}{240,255,240}
\definecolor{box-green-border}{RGB}{60,150,60}
\definecolor{box-yellow}{RGB}{255,253,230}
\definecolor{box-yellow-border}{RGB}{180,150,0}

% ─── Listings ─────────────────────────────────────────────────────────────────
\lstdefinelanguage{DSL}{
  keywords={def-table, libro, hoja, tabla, load, list, defparameter,
            conditional-rendering, exists, nav, ultima-fila, primera-fila,
            xl-generate, xl-run-generated},
  morekeywords=[2]{countif, counta, sum-range, str, col, trange,
                   _equals, _and, _or},
  keywordstyle=\color{lisp-kw}\bfseries,
  keywordstyle=[2]\color{lisp-mac},
  stringstyle=\color{lisp-str},
  commentstyle=\color{lisp-cmt}\itshape,
  morestring=[b]",
  morecomment=[l]{;;},
  sensitive=false,
  basicstyle=\ttfamily\small,
  backgroundcolor=\color{lisp-bg},
  frame=single,
  framerule=0.4pt,
  rulecolor=\color{gray!40},
  breaklines=true,
  numbers=left,
  numberstyle=\tiny\color{gray},
  numbersep=8pt,
  xleftmargin=20pt,
  xrightmargin=4pt,
  aboveskip=6pt,
  belowskip=6pt,
  showstringspaces=false,
  tabsize=2,
}

\lstset{language=DSL}

% ─── Cajas ────────────────────────────────────────────────────────────────────
\tcbuselibrary{skins,breakable}

\newtcolorbox{infobox}[1]{
  colback=box-blue, colframe=box-blue-border,
  fonttitle=\bfseries, title=#1,
  breakable, sharp corners, left=6pt, right=6pt, top=4pt, bottom=4pt
}

\newtcolorbox{decisionbox}[1]{
  colback=box-yellow, colframe=box-yellow-border,
  fonttitle=\bfseries, title=#1,
  breakable, sharp corners, left=6pt, right=6pt, top=4pt, bottom=4pt
}

\newtcolorbox{sintaxisbox}[1]{
  colback=box-green, colframe=box-green-border,
  fonttitle=\bfseries\ttfamily, title=Sintaxis: \texttt{#1},
  breakable, sharp corners, left=6pt, right=6pt, top=4pt, bottom=4pt
}

% ─── Títulos ──────────────────────────────────────────────────────────────────
\titleformat{\chapter}[hang]{\large\bfseries}{\thechapter.}{0.5em}{}[\vspace{2pt}\hrule]
\titleformat{\section}[hang]{\normalsize\bfseries}{\thesection}{0.5em}{}
\titleformat{\subsection}[hang]{\normalsize\itshape\bfseries}{\thesubsection}{0.5em}{}

\hypersetup{hidelinks}
\setlength{\parskip}{6pt}

% ──────────────────────────────────────────────────────────────────────────────
\begin{document}

\chapter{Tutorial: construcción de un libro de cálculo con el DSL}
\label{cap:tutorial}

Este cap\'{\i}tulo muestra, paso a paso, c\'omo construir un programa que
genere autom\'aticamente un libro de c\'alculo a partir de la descripci\'on
de un problema de horario.
Al final del tutorial tendr\'as el archivo \texttt{Horario\_Tesis.xlsx}
en disco, listo para abrir.

% ══════════════════════════════════════════════════════════════════════════════
\section{Instalaci\'on}
\label{sec:requisitos}

Instala los dos programas siguientes.

\textbf{Steel Bank Common Lisp (SBCL)} --- descarga desde \url{https://www.sbcl.org/} e instala seg\'un las instrucciones de tu sistema operativo.

\textbf{Python 3 con openpyxl}:

\begin{lstlisting}[language=bash, numbers=none, backgroundcolor=\color{gray!10},
                   basicstyle=\ttfamily\small, keywordstyle={}]
pip install openpyxl
\end{lstlisting}

Descarga el generador.
Los archivos que no debes modificar son la biblioteca del sistema.
Tus dos archivos van en la \textbf{misma carpeta} que la biblioteca:

\begin{lstlisting}[language={}, numbers=none, backgroundcolor=\color{gray!10},
                   basicstyle=\ttfamily\small]
generador/
+-- dsl-directo.lisp            <- biblioteca (no modificar)
+-- ast-def.lisp                <- biblioteca (no modificar)
+-- generate-code-direct.lisp   <- biblioteca (no modificar)
+-- code-generation-utils.lisp  <- biblioteca (no modificar)
+-- hoja_con_formulas.py        <- biblioteca (no modificar)
+-- data-tutorial.lisp          <- tus datos  (editar)
\-- tutorial-horario.lisp       <- tu programa (editar)
\end{lstlisting}

No crees subcarpetas.
El programa se ejecuta as\'{\i} desde esa misma carpeta:

\begin{lstlisting}[language=bash, numbers=none, backgroundcolor=\color{gray!10},
                   basicstyle=\ttfamily\small, keywordstyle={}]
cd generador/
sbcl --load tutorial-horario.lisp
\end{lstlisting}

% ══════════════════════════════════════════════════════════════════════════════
\section{El problema}
\label{sec:problema}

La Facultad de Matem\'atica y Computaci\'on de la Universidad de La Habana
celebra sus defensas de tesis en junio de 2026.
Se tienen 13 defensas programadas en distintos locales (Postgrado, Sal\'on
decanato, Francofon\'{\i}a) durante los d\'{\i}as 5, 8 y 10 de junio.
Cada defensa tiene un estudiante, un tribunal de cinco profesores
(tutor, presidente, secretario, vocal y oponente), un d\'{\i}a, una hora y un local.
Participan 32 profesores en total, algunos en m\'as de una defensa.
Se quiere saber cu\'antas defensas acumula cada profesor
y detectar autom\'aticamente si alg\'un profesor integra dos tribunales
programados en el mismo d\'{\i}a y hora.

El libro tendr\'a dos tablas en la misma hoja:
una con las defensas y otra con los profesores.

% ══════════════════════════════════════════════════════════════════════════════
\section{Construyendo el programa paso a paso}
\label{sec:pasos}

Crea el archivo \texttt{tutorial-horario.lisp} en la carpeta del generador
y ve a\~nadiendo cada bloque seg\'un se describe a continuaci\'on.

% ─────────────────────────────────────────────────────────────────────────────
\subsection{Paso 1 --- cargar el DSL}

La primera l\'{\i}nea del fichero carga el sistema:

\begin{lstlisting}[caption={tutorial-horario.lisp --- inicio.}]
(load "dsl-directo.lisp")
\end{lstlisting}

\texttt{load} ejecuta el archivo indicado.
\texttt{dsl-directo.lisp} carga a su vez el resto de la biblioteca
(ast-def, generate-code-direct, etc.) y deja disponibles todas las formas
del DSL: \texttt{def-table}, \texttt{libro}, \texttt{xl-generate}\ldots

% ─────────────────────────────────────────────────────────────────────────────
\subsection{Paso 2 --- preparar los datos}

Crea el archivo \texttt{data-tutorial.lisp} en la misma carpeta.
\texttt{defparameter} define una variable global con los datos de entrada.
El formato es una lista de listas: la fila 1 y 2 son prefijos vac\'{\i}os
que el generador coloca antes de los datos, la fila 3 son los encabezados
de columna, y a partir de la fila 4 van los datos reales.
Las celdas que el generador va a calcular autom\'aticamente se dejan como
cadena vac\'{\i}a \texttt{""} ---el sistema las sobreescribir\'a con f\'ormulas.

\begin{lstlisting}[caption={data-tutorial.lisp (extracto).}]
(defparameter *DATOS-DEFENSAS*
  '(("" "" "" "" "" "" "" "" "")
    ("" "" "" "" "" "" "" "" "")
    ("Dia" "Hora" "Estudiante" "Tutor"
     "Presidente" "Secretario" "Vocal" "Oponente" "Local")
    ;; --- 5 de junio ---
    ("05/06/2026" "10:00"
     "Angel Daniel Alonso Guevara"
     "MSc. Carmen T. Fernandez Montoto"
     "Dr. Alberto Fernandez Oliva"
     "Lic. Alejandro Labourdette-Lantigua Soto"
     "Lic. Dennis Daniel Gonzalez Duran"
     "MSc. Aracelys Garcia Armenteros"
     "Postgrado")
    ;; --- 8 de junio ---
    ("08/06/2026" "09:30"
     "Adrian Hernandez Castellanos"
     "Lic. Alejandra Monzon Pena"
     "Lic. Amanda Noris Hernandez"
     "Lic. Kevin Manzano Rodriguez"
     "Lic. Daniel Abad Fundora"
     "Lic. Rodrigo Garcia Gomez"
     "Salon decanato")
    ... ;; 11 filas mas en data-tutorial.lisp
    ))

(defparameter *DATOS-PROFESORES*
  '(("" "")
    ("" "")
    ("Nombre" "Defensas")
    ("Dr. Alberto Fernandez Oliva"          "")
    ("DrC. Alejandro Piad Morfis"           "")
    ("Lic. Alejandra Monzon Pena"           "")
    ("MSc. Fernando Raul Rodriguez Flores"  "")
    ... ;; 28 profesores mas en data-tutorial.lisp
    ))
\end{lstlisting}

El archivo completo con las 13 defensas y los 32 profesores est\'a en
\texttt{data-tutorial.lisp}.
El nombre de cada profesor incluye su grado acad\'emico
(\texttt{MSc.}, \texttt{Lic.}, \texttt{DrC.}, etc.) para que la
b\'usqueda que genera el sistema lo encuentre exactamente en las columnas
de tribunal, donde los nombres tambi\'en llevan grado.

A\~nade en \texttt{tutorial-horario.lisp} la l\'{\i}nea que carga los datos:

\begin{lstlisting}[caption={tutorial-horario.lisp --- cargar datos.},
                   firstnumber=2]
(load "data-tutorial.lisp")
\end{lstlisting}

A partir de este momento tienes disponibles las variables
\texttt{*DATOS-DEFENSAS*} y \texttt{*DATOS-PROFESORES*}.
Si a\~nadieras una defensa m\'as, bastar\'{\i}a con a\~nadir una fila
a \texttt{*DATOS-DEFENSAS*} ---el resto del programa no cambia.

% ─────────────────────────────────────────────────────────────────────────────
\subsection{Paso 3 --- declarar la tabla de defensas}

Cada defensa es una fila.
Las columnas siguen el orden del Excel fuente.
Los cinco roles de tribunal van juntos y consecutivos para poder
referenciarlos con un solo rango en las f\'ormulas.

\begin{lstlisting}[caption={tutorial-horario.lisp --- tabla de defensas.},
                   firstnumber=4]
(def-table defensas-table
  ((dia "") (hora "") (estudiante "") (tutor "")
   (presidente "") (secretario "") (vocal "") (oponente "") (local ""))
  :cell-height 1
  :cell-width  1)
\end{lstlisting}

\textbf{Qu\'e hace \texttt{def-table}:}
Cada par \texttt{(nombre "")} declara una columna.
El primer elemento (\texttt{dia}, \texttt{tutor}...) es el \emph{nombre
interno}: el s\'{\i}mbolo con el que te referir\'as a esa columna en
cualquier f\'ormula del resto del programa.
El segundo elemento (\texttt{""}) es el valor por defecto de la celda;
en esta tabla todos son vac\'{\i}os porque los datos vienen de
\texttt{*DATOS-DEFENSAS*}.

Si quisieras a\~nadir una columna \texttt{observacion}, bastar\'{\i}a con
incluir \texttt{(observacion "")} en la lista y a\~nadir el valor
correspondiente en cada fila de \texttt{*DATOS-DEFENSAS*}.

% ─────────────────────────────────────────────────────────────────────────────
\subsection{Paso 4 --- declarar la tabla de profesores}

Cada profesor es una fila.
La columna \texttt{defensas} se calcula autom\'aticamente.
El campo \texttt{nombre} almacena el nombre completo con grado
(\texttt{"MSc.\ Fernando Ra\'ul Rodr\'{\i}guez Flores"}) para que
coincida exactamente con los valores en las columnas de tribunal de
\texttt{defensas-table}.

\begin{lstlisting}[caption={tutorial-horario.lisp --- tabla de profesores.},
                   firstnumber=10]
(def-table profesores-table
  ((nombre "") (defensas ""))
  :cell-height 1
  :cell-width  1
  :computed
    ((defensas (countif (trange defensas-table tutor oponente)
                        (col nombre))))
  :render
    (conditional-rendering
      :condition
        (exists (r :rows-of defensas-table)
          :self-in  (tutor presidente secretario vocal oponente)
          :matching (dia hora))
      :target-columns (nombre)))
\end{lstlisting}

\textbf{Qu\'e hace \texttt{:computed}:}
Declara columnas cuyo valor es una f\'ormula, no un dato manual.
La forma \texttt{(countif rango criterio)} genera una f\'ormula
\texttt{COUNTIF} en Excel.
\texttt{(trange defensas-table tutor oponente)} produce el rango que
abarca desde la columna \texttt{tutor} hasta la columna \texttt{oponente}
en todas las filas de \texttt{defensas-table}.
\texttt{(col nombre)} devuelve el valor de la columna \texttt{nombre}
\emph{en la fila actual} de \texttt{profesores-table}.
Combinados, la f\'ormula cuenta cu\'antas veces aparece este profesor
en cualquiera de los cinco roles.

\textbf{Modificaci\'on: contar por rol.}
Si quisieras saber adem\'as cu\'antas veces el profesor fue tutor
espec\'{\i}ficamente, a\~nadir\'{\i}as una columna \texttt{como-tutor}:

\begin{lstlisting}[caption={Columna adicional: solo tutores.},numbers=none]
(def-table profesores-table
  ((nombre "") (defensas "") (como-tutor ""))
  ...
  :computed
    ((defensas    (countif (trange defensas-table tutor oponente)
                           (col nombre)))
     (como-tutor  (countif (trange defensas-table tutor tutor)
                           (col nombre)))))
\end{lstlisting}

El segundo \texttt{countif} usa \texttt{(trange defensas-table tutor tutor)}
---un rango de una sola columna--- para contar solo apariciones en la columna
\texttt{tutor}.
Para contar como presidente usar\'{\i}as \texttt{(trange ... presidente
presidente)}, y as\'{\i} con cada rol.

\textbf{Qu\'e hace \texttt{:render}:}
Asocia una condici\'on de coloreado con una o m\'as columnas destino.
La condici\'on usa \texttt{exists}: busca si existe alguna fila en
\texttt{defensas-table} donde \texttt{:self-in} ---los roles--- contienen
el nombre del profesor, y cuyo slot \texttt{dia + hora} aparece en m\'as
de una defensa simult\'anea (\texttt{:matching}).
El sistema genera autom\'aticamente una \texttt{FormulaRule} en el
\texttt{.xlsx} que se reevalúa cada vez que el usuario edita el libro.

% ─────────────────────────────────────────────────────────────────────────────
\subsection{Paso 5 --- ensamblar el libro}

Las dos tablas van en la misma hoja, una al lado de la otra.
Debajo de cada tabla se colocan etiquetas y totales.

\begin{lstlisting}[caption={tutorial-horario.lisp --- workbook.},
                   firstnumber=25]
(libro horario-tesis
  :filename "Horario_Tesis.xlsx"
  :hojas (list
    (hoja "Defensas"
      (tabla defensas-table   :data *DATOS-DEFENSAS*)
      (tabla profesores-table :data *DATOS-PROFESORES*)
      :fixed-expressions
        (((nav (ultima-fila defensas-table dia) abajo 1)
          (str "Total defensas:"))
         ((nav (ultima-fila defensas-table dia) abajo 1 derecha 2)
          (counta (trange defensas-table estudiante)))
         ((nav (ultima-fila profesores-table nombre) abajo 1)
          (str "Total profesores:"))
         ((nav (ultima-fila profesores-table nombre) abajo 1 derecha 1)
          (sum-range (trange profesores-table defensas)))))))
\end{lstlisting}

\textbf{Qu\'e hace cada forma:}
\texttt{libro} crea el workbook con un nombre de archivo de salida.
\texttt{hoja} define una pesta\~na; puede recibir m\'as de una
\texttt{tabla}, en cuyo caso las coloca en la misma hoja una al lado de otra.
\texttt{tabla} instancia una plantilla (\texttt{def-table}) con datos
concretos.

\texttt{:fixed-expressions} coloca f\'ormulas o etiquetas en celdas fuera
de las tablas.
La posici\'on se expresa con \texttt{nav} y \texttt{ultima-fila}:
\texttt{(nav (ultima-fila defensas-table dia) abajo 1)} calcula la celda
que est\'a exactamente una fila debajo de la \'ultima fila de datos de
\texttt{defensas-table}, columna \texttt{dia}.
\texttt{derecha 2} a\~nade un desplazamiento de dos columnas a la derecha.
Estas posiciones se recalculan autom\'aticamente si el n\'umero de filas
cambia, as\'{\i} que si a\~nades m\'as defensas los totales se desplazan solos.

% ─────────────────────────────────────────────────────────────────────────────
\subsection{Paso 6 --- generar y ejecutar}

A\~nade las dos \'ultimas l\'{\i}neas al fichero:

\begin{lstlisting}[caption={tutorial-horario.lisp --- generar y ejecutar.},
                   firstnumber=44]
(xl-generate horario-tesis "horario-tesis.py")
(xl-run-generated "horario-tesis.py")
\end{lstlisting}

\texttt{xl-generate} recorre el \'arbol AST que construyeron las macros
anteriores y produce el script Python \texttt{horario-tesis.py}.
\texttt{xl-run-generated} ejecuta ese script con Python 3 y escribe
el archivo \texttt{Horario\_Tesis.xlsx} en disco.

Ejecuta desde la carpeta del generador (la misma donde est\'an todos
los archivos \texttt{.lisp}):

\begin{lstlisting}[language=bash, numbers=none, backgroundcolor=\color{gray!10},
                   basicstyle=\ttfamily\small, keywordstyle={}]
sbcl --load tutorial-horario.lisp
\end{lstlisting}

% ══════════════════════════════════════════════════════════════════════════════
\section{Resultado}
\label{sec:resultado}

El archivo \texttt{Horario\_Tesis.xlsx} aparece en la carpeta
\texttt{mi-horario/}.
\'Abrelo con Excel o LibreOffice Calc.

\begin{center}
\renewcommand{\arraystretch}{1.3}
\begin{tabular}{ll}
\toprule
\textbf{Rango} & \textbf{Contenido} \\
\midrule
Columnas A--I, filas 4--16 & \texttt{defensas-table}: 13 defensas de junio \\
Columna J                  & Separador en blanco \\
Columnas K--L, filas 4--35 & \texttt{profesores-table}: 32 profesores \\
Columna L, filas 4--35     & F\'ormulas \texttt{COUNTIF} generadas \\
Fila 17, columnas A y C    & Total de defensas con \texttt{COUNTA} \\
Fila 37, columnas K--L     & Total de profesores con \texttt{SUM} \\
Rango K4:K35               & \texttt{FormulaRule}: nombres en rojo si hay conflicto \\
\bottomrule
\end{tabular}
\end{center}

Con este horario real, cada profesor aparece en entre 1 y 4 defensas
(\texttt{MSc.\ Fernando Ra\'ul Rodr\'{\i}guez Flores} y
\texttt{MSc.\ Aracelys Garc\'{\i}a Armenteros} participan en 4 cada uno).
Los slots \texttt{10/06/2026 09:30} y \texttt{10/06/2026 13:30} tienen
dos defensas simult\'aneas en distintos locales, pero los tribunales son
completamente distintos, por lo que ning\'un profesor aparece en rojo:
el sistema verifica correctamente que no hay conflictos en este horario.

% ══════════════════════════════════════════════════════════════════════════════
\section{El fichero completo}
\label{sec:completo}

\begin{lstlisting}[caption={tutorial-horario.lisp completo.}]
(load "dsl-directo.lisp")
(load "data-tutorial.lisp")

;; Tabla 1: una fila por defensa
;; Columnas en el orden del Excel fuente.
;; Los roles de tribunal son consecutivos (tutor...oponente) para
;; referenciarlos con un unico trange en el COUNTIF y en el exists.
(def-table defensas-table
  ((dia "") (hora "") (estudiante "") (tutor "")
   (presidente "") (secretario "") (vocal "") (oponente "") (local ""))
  :cell-height 1
  :cell-width  1)

;; Tabla 2: una fila por profesor
;; El nombre incluye el grado ("MSc. ...", "Lic. ...") para que el
;; COUNTIF lo encuentre tal cual en las columnas de tribunal.
(def-table profesores-table
  ((nombre "") (defensas ""))
  :cell-height 1
  :cell-width  1
  :computed
    ((defensas (countif (trange defensas-table tutor oponente)
                        (col nombre))))
  :render
    (conditional-rendering
      :condition
        (exists (r :rows-of defensas-table)
          :self-in  (tutor presidente secretario vocal oponente)
          :matching (dia hora))
      :target-columns (nombre)))

(libro horario-tesis
  :filename "Horario_Tesis.xlsx"
  :hojas (list
    (hoja "Defensas"
      (tabla defensas-table   :data *DATOS-DEFENSAS*)
      (tabla profesores-table :data *DATOS-PROFESORES*)
      :fixed-expressions
        (((nav (ultima-fila defensas-table dia) abajo 1)
          (str "Total defensas:"))
         ((nav (ultima-fila defensas-table dia) abajo 1 derecha 2)
          (counta (trange defensas-table estudiante)))
         ((nav (ultima-fila profesores-table nombre) abajo 1)
          (str "Total profesores:"))
         ((nav (ultima-fila profesores-table nombre) abajo 1 derecha 1)
          (sum-range (trange profesores-table defensas)))))))

(xl-generate horario-tesis "horario-tesis.py")
(xl-run-generated "horario-tesis.py")
\end{lstlisting}

% ══════════════════════════════════════════════════════════════════════════════
%  A partir de aqu\'i: manual de referencia
% ══════════════════════════════════════════════════════════════════════════════

\section{C\'omo funciona el sistema}
\label{sec:conceptos}

Al cargar \texttt{tutorial-horario.lisp}, SBCL ejecuta las macros del DSL
y construye en memoria un \textbf{\'arbol de sintaxis abstracta (AST)}.
Cada nodo del \'arbol representa un concepto del problema: una tabla, una
expresi\'on de celda, una regla de coloreado.

\texttt{xl-generate} recorre ese \'arbol y produce un script Python.
\texttt{xl-run-generated} ejecuta el script, que usa la biblioteca
\texttt{openpyxl} para escribir el archivo \texttt{.xlsx}.

El DSL no contiene ninguna referencia al formato de salida: el mismo programa
podr\'{\i}a alimentar un generador diferente (PDF, HTML, LibreOffice) sin
ning\'un cambio en el c\'odigo del usuario.

% ══════════════════════════════════════════════════════════════════════════════
\section{Gu\'{\i}a de las formas del DSL}
\label{sec:guia}

% ─────────────────────────────────────────────────────────────────────────────
\subsection{\texttt{def-table} --- definir una tabla}

\begin{sintaxisbox}{def-table}
\begin{lstlisting}
(def-table nombre-tabla
  ((col1 "Cabecera1") (col2 "Cabecera2") ...)
  :cell-height N          ; filas fisicas por fila logica (defecto 1)
  :cell-width  N          ; columnas fisicas por columna logica (defecto 1)
  :computed               ; columnas calculadas (opcional)
    ((col-x (expresion)))
  :render                 ; regla de coloreado (opcional)
    (conditional-rendering ...))
\end{lstlisting}
\end{sintaxisbox}

Cada par \texttt{(columna "")} declara el nombre interno del DSL y el texto
del encabezado visual.
La clave \texttt{:computed} asocia columnas a expresiones que el generador
traduce a f\'ormulas.
La clave \texttt{:render} acepta una forma \texttt{conditional-rendering}
que especifica bajo qu\'e condici\'on se colorean determinadas columnas.

% ─────────────────────────────────────────────────────────────────────────────
\subsection{\texttt{exists} --- condici\'on de existencia}

\begin{sintaxisbox}{exists :rows-of}
\begin{lstlisting}
(exists (var :rows-of tabla-id)
  :self-in  (col1 col2 ...)    ; el valor actual aparece en alguna de estas
  :matching (clave1 clave2 ...)); ese slot tiene mas de una fila programada
\end{lstlisting}
\end{sintaxisbox}

\begin{sintaxisbox}{exists :all-rows}
\begin{lstlisting}
(exists (var :all-rows)
  :matching    (col1 col2 ...)  ; columnas que deben coincidir entre filas
  :any-overlap (col1 col2 ...)) ; columnas donde algun valor debe compartirse
\end{lstlisting}
\end{sintaxisbox}

\texttt{:rows-of} busca en otra tabla; \texttt{:all-rows} busca en la misma tabla.
El generador produce una f\'ormula \texttt{SUMPRODUCT} que persiste en el
archivo y se reevalúa autom\'aticamente cuando el usuario edita los datos.

% ─────────────────────────────────────────────────────────────────────────────
\subsection{\texttt{nav} --- posicionamiento relativo}

\begin{sintaxisbox}{nav / ultima-fila / primera-fila}
\begin{lstlisting}
(nav (ultima-fila  tabla-id columna) abajo  N)
(nav (ultima-fila  tabla-id columna) abajo  N derecha M)
(nav (primera-fila tabla-id columna) arriba N)
\end{lstlisting}
\end{sintaxisbox}

Calcula la posici\'on de una celda como un ancla en una tabla m\'as
pasos direccionales (\texttt{abajo}, \texttt{arriba}, \texttt{derecha},
\texttt{izquierda}).
Las coordenadas se recalculan cada vez que cambia el n\'umero de filas
de datos, por lo que los totales de pie de tabla se desplazan solos.

% ─────────────────────────────────────────────────────────────────────────────
\subsection{\texttt{libro} / \texttt{hoja} / \texttt{tabla} --- ensamblar el workbook}

\begin{sintaxisbox}{libro / hoja / tabla}
\begin{lstlisting}
(libro nombre-workbook
  :filename "nombre.xlsx"
  :hojas (list
    (hoja "Nombre hoja"
      (tabla plantilla1 :data *DATOS-1*)
      (tabla plantilla2 :data *DATOS-2*)
      :fixed-expressions
        (((nav ancla dir N) (expresion))
         ...))))
\end{lstlisting}
\end{sintaxisbox}

Cuando una \texttt{hoja} recibe m\'as de una \texttt{tabla}, el generador
las coloca lado a lado separadas por una columna en blanco.
\texttt{:fixed-expressions} coloca expresiones fuera de las tablas en
posiciones relativas calculadas con \texttt{nav}.

% ══════════════════════════════════════════════════════════════════════════════
\section{Referencia r\'apida de expresiones}
\label{sec:referencia}

\begin{center}
\renewcommand{\arraystretch}{1.4}
\begin{tabular}{>{\ttfamily}p{5.6cm}p{7.4cm}}
\toprule
\textbf{Forma} & \textbf{Descripci\'on} \\
\midrule
(col nombre)              & Valor de la columna en la fila actual \\
(trange tabla desde hasta)& Rango en otra tabla \\
(countif rango criterio)  & F\'ormula COUNTIF \\
(counta rango)            & F\'ormula COUNTA \\
(sum-range rango)         & F\'ormula SUM \\
(str "texto")             & Literal de texto en celda \\
(\_if cond entonces sino) & Condicional de tres ramas \\
(non-empty expr)          & Verdadero si la expresi\'on no est\'a vac\'{\i}a \\
(show-nothing)            & Produce la cadena vac\'{\i}a \texttt{""} \\
(equals a b)              & Igualdad \\
(different a b)           & Desigualdad \\
(\_and a b) / (\_or a b)  & Conjunci\'on y disyunci\'on \\
(add a b) / (subtract a b)& Suma y resta \\
(multiply a b) / (divide a b) & Multiplicaci\'on y divisi\'on \\
(time-add hora minutos)   & Suma minutos a una hora \\
(concat a b)              & Concatenaci\'on de texto \\
(previous-row expr)       & Valor en la fila anterior \\
(next-row expr)           & Valor en la fila siguiente \\
(first-row)               & Verdadero en la primera fila de datos \\
(tcol tabla nombre)       & Columna de una tabla concreta (desambiguaci\'on) \\
(lookup clave campo)      & B\'usqueda estilo VLOOKUP \\
(xl-generate wb archivo)  & Genera el script de construcci\'on \\
(xl-run-generated archivo)& Ejecuta el script y crea el \texttt{.xlsx} \\
\bottomrule
\end{tabular}
\end{center}

% ══════════════════════════════════════════════════════════════════════════════
\section{Formas adicionales del DSL}
\label{sec:macros-extra}

% ─────────────────────────────────────────────────────────────────
\subsection*{Plantillas param\'etricas}

Cuando la misma estructura de tabla se necesita varias veces con un detalle
variable (por ejemplo, una tabla de franjas por cada d\'{\i}a de la semana),
se usan \texttt{:inst-params}:

\begin{lstlisting}
(def-table franja-dia
  ((turno "") (salon "") (grupo ""))
  :inst-params (dia))

(hoja "Semana"
  (tabla franja-dia :dia 'lunes  :data *datos-lunes*)
  (tabla franja-dia :dia 'martes :data *datos-martes*))
\end{lstlisting}

% ─────────────────────────────────────────────────────────────────
\subsection*{Celdas de varias filas o columnas}

\texttt{:cell-height~N} hace que cada fila l\'ogica ocupe \texttt{N} filas
f\'{\i}sicas en la hoja generada.

\begin{lstlisting}
(def-table horario-visual
  ((hora "") (actividad "") (sala ""))
  :cell-height 2)
\end{lstlisting}

% ─────────────────────────────────────────────────────────────────
\subsection*{Referencias cruzadas entre hojas}

\begin{center}
\begin{tabular}{lp{8cm}}
\toprule
\textbf{Forma} & \textbf{Para qu\'e sirve} \\
\midrule
\texttt{(sheet-ref hoja plantilla-celda)} & Referencia a una celda de
  otra hoja con nombre fijo. \\
\texttt{(cross-cell sheet xcol row)} & Referencia a una celda de una hoja
  enlazada din\'amicamente mediante \texttt{collect-over}. \\
\bottomrule
\end{tabular}
\end{center}

\end{document}
